% =============================================================================
% SECTION 2: RELATED WORK
% =============================================================================

\section{Related Work}
\label{sec:related}

This section organizes the literature around four competing approaches to the same
problem: generating text that is both fluent and factually grounded. Each subsection
identifies what the approach achieves and where it stops, building toward the specific
gap that DCA-Trie addresses.

% -----------------------------------------------------------------------------
\subsection{Hallucination Mitigation}
\label{sec:related:hallucination}
% -----------------------------------------------------------------------------

The most widely deployed defense against hallucination is retrieval-augmented generation
(RAG)~\citep{lewis2020rag}. A dense retriever selects passages by embedding similarity,
and the LLM generates from the retrieved context. Extensions such as
FLARE~\citep{jiang2023flare} trigger retrieval when next-token confidence drops, and
IRCoT~\citep{trivedi2022ircot} interleaves retrieval with chain-of-thought steps to
improve multi-hop accuracy. Self-RAG~\citep{asai2024selfrag} adds reflection tokens
that control when to retrieve and whether retrieved passages support the generated
claims.

These methods share a structural limitation. The retriever determines what the LLM
\textit{sees}; it does not determine what the LLM is \textit{permitted to generate}.
Because the generation mechanism remains unchanged, any token in the vocabulary retains
a positive probability at every step, including tokens that correspond to no verified
fact. Structural faithfulness with probability one cannot be achieved under any
architecture that modifies only the input context without modifying the token selection
mechanism~\citep{geng2023grammarconstrained, luo2025-graph-constrained-reasoning}.

Chain-of-thought (CoT) prompting~\citep{wei2022cot} takes a different approach:
rather than retrieving facts, it elicits intermediate reasoning steps that decompose
complex questions into simpler sub-problems. Self-consistency~\citep{wang2023selfconsistency}
samples multiple independent CoT chains and selects the final answer by majority vote,
reducing sensitivity to individual hallucinated chains. Tree-of-Thought~\citep{yao2023-tree-of-thought}
treats generation as search through a tree of partial reasoning sequences, using the LLM
as both generator and evaluator for pruning.

All prompting-based methods modify the \textit{input} to the LLM without modifying the
\textit{generation mechanism}. A CoT chain can be entirely fluent and internally coherent
while containing fabricated intermediate steps. Self-consistency reduces this risk through
redundancy, but a majority of independently hallucinated chains can still produce an
incorrect answer. Tree-of-Thought relies on the model's own parameters for pruning
decisions, without external verification. True structural faithfulness requires that
invalid tokens receive exactly zero probability, achievable only by altering the decoding
algorithm itself~\citep{geng2023grammarconstrained, willard2023outlines}.

% -----------------------------------------------------------------------------
\subsection{Knowledge Graph Reasoning}
\label{sec:related:kg}
% -----------------------------------------------------------------------------

Knowledge graphs encode verified facts as directed, labelled triplets
$(e, r, e')$~\citep{bollacker2008freebase}. KGQA requires identifying the answer
entity by traversing chains of these triplets. Benchmarks such as
WebQSP~\citep{yih2016websqp} and Complex WebQuestions~\citep{talmor2018cwq} evaluate
multi-hop path traversal over Freebase, demanding adherence to graph topology.

Early KGQA methods relied on semantic parsing: transforming questions into logical
queries (e.g., SPARQL) for execution over the
KG~\citep{lan2022survey}. These approaches guarantee faithfulness because answers are
produced by verified graph operations. The requirement of training on ground-truth
logical forms, and the non-executability of model-generated queries due to syntax or
semantic errors, limits their practical deployment~\citep{lan2022survey}.

Retrieval-based methods use LLMs to extract relevant subgraphs or paths from the KG.
RoG~\citep{luo2024rog} trains a reasoning path planner that retrieves candidate paths
and generates natural language explanations. GNN-RAG~\citep{mavromatis2022rearev}
combines graph neural networks with language models for retrieval over structured data.
These methods improve relevance by filtering the KG, but the retrieved context remains
a soft signal: it influences generation without constraining it.

Agent-based methods treat the LLM as a reasoning agent that iteratively queries the
KG. ToG~\citep{sun2024tog} prompts the LLM to explore relevant facts hop-by-hop using
beam search on the KG. StructGPT~\citep{jiang2023structgpt} interfaces the LLM with
structured data through iterative retrieval. These approaches are flexible but depend
on the LLM's instruction-following ability and are slow due to multiple KG queries per
question.

The shared limitation across retrieval-based and agent-based KG reasoning is that
retrieved context is a soft influence. The retriever or agent selects what the LLM
observes; the LLM retains the freedom to generate any token. Even with perfect retrieval,
hallucination remains structurally possible.

% -----------------------------------------------------------------------------
\subsection{Constrained Decoding}
\label{sec:related:constrained}
% -----------------------------------------------------------------------------

Constrained decoding restricts generation to a predefined valid set at each step by
directly intervening in the decoding mechanism. Unlike prompting or retrieval, which
modify the input to the LLM, constrained decoding alters the token selection process
itself. Invalid tokens receive zero probability before sampling, making their generation
structurally impossible.

\textbf{Format-constrained methods.} Grammar-constrained decoding
(GCD)~\citep{geng2023grammarconstrained} constrains generation to strings conforming to
a context-free grammar via an Earley parser oracle. Outlines~\citep{willard2023outlines}
reduces the computational cost with a finite-state machine (FSM) formulation,
precomputing an index from automaton states to valid vocabulary subsets for amortised
$O(1)$ lookup. Additional methods include Synchromesh~\citep{poesia2022synchromesh} for
program synthesis, PICARD~\citep{scholak2021picard} for schema-constrained SQL parsing,
and NeuroLogic Decoding~\citep{lu2021neurologic} for propositional logic constraints.

These methods enforce output \textit{format} rather than \textit{factual content}. A
grammatically valid SQL query can reference non-existent tables. A well-formed entity
name can be factually incorrect. The extension from format constraints to KG path
constraints requires a constraint oracle that encodes the relational structure of the
knowledge graph itself.

\textbf{KG-constrained methods.} Graph-Constrained Reasoning
(GCR)~\citep{luo2025-graph-constrained-reasoning} builds a KG-Trie: a prefix tree of
all KG paths within $L$ hops of the question entities, created by breadth-first search
before generation. The trie is used as the constraint oracle during beam-search
decoding. GCR achieves 92.6\% Hits@1 on WebQSP with 100\% structural faithfulness
using a Llama-3.1-8B backbone. The trie is constructed once and never updated: the
valid token set is fixed regardless of what has been generated.

Decoding on Graphs (DoG)~\citep{li2024-decoding-on-graphs} addresses GCR's static
oracle through step-wise trie expansion. After locking in an entity, the constraint
trie is expanded with all KG neighbors of that entity, making the oracle topologically
dynamic. DoG achieves 91.0\% Hits@1 on WebQSP with 100\% structural faithfulness.
However, the expansion rule admits all neighbors regardless of relevance to the question,
and step-wise trie reconstruction introduces computational overhead.

RouterKGQA~\citep{yuan2026routerkgqa} routes between a specialized KG-reasoning model
and a general LLM based on question complexity, using SBERT embeddings for post-hoc
answer filtering. RouterKGQA's semantic scoring operates at the answer selection level
rather than at the token constraint level: paths are generated first, then filtered
afterward.

\textbf{The gap.} All three KG-constrained methods determine their constraint sets
from graph structure and question entities alone. None conditions the valid token set on
question semantics or the partially generated chain. The constraint set is indifferent
to whether a path is actually related to the question being asked. On multi-hop questions,
this produces valid token sets containing thousands of structurally correct but
semantically irrelevant paths, forcing the LLM to navigate among them using its own
knowledge and reintroducing the hallucination risk that constrained decoding was designed
to eliminate.

% -----------------------------------------------------------------------------
\subsection{Comparison and Research Gap}
\label{sec:related:gap}
% -----------------------------------------------------------------------------

Table~\ref{tab:related_comparison} compares the approaches reviewed above along four
dimensions: whether the method enforces structural faithfulness (the generation
mechanism, not just the input), whether the constraint set adapts during decoding,
whether semantic relevance enters the constraint formation, and the type of oracle used.

\begin{table}[htpb]
  \centering
  \caption{Comparison of KG-grounded reasoning methods. \textbf{Dynamic} indicates
  whether the constraint set changes during decoding. \textbf{Semantic} indicates
  whether question meaning enters the constraint oracle. The oracle type shows what
  drives constraint formation.}
  \label{tab:related_comparison}
  \small
  \begin{tabular}{@{}lcccl@{}}
    \toprule
    Method        & Structural & Dynamic & Semantic & Oracle Type    \\
    \midrule
    RAG           & No         & No      & Soft     & Retriever      \\
    CoT           & No         & No      & No       & None           \\
    Self-RAG      & No         & No      & Soft     & Reflection     \\
    GCD           & Format     & No      & No       & Grammar/FSM    \\
    Outlines      & Format     & No      & No       & FSM-index      \\
    GCR           & Yes        & No      & No       & Static KG-Trie \\
    DoG           & Yes        & Partial & No       & Topology       \\
    RouterKGQA    & Prob.      & No      & Yes      & SBERT filter   \\
    \textbf{DCA-Trie} & \textbf{Yes} & \textbf{Yes} & \textbf{Yes} & \textbf{TypeOracle} \\
    \bottomrule
  \end{tabular}
\end{table}

The table reveals a structural division. Methods above GCR achieve either semantic
awareness (RAG, RouterKGQA) or structural faithfulness (GCD, Outlines), but not both.
GCR and DoG achieve structural faithfulness but lack semantic awareness: their constraint
oracles are indifferent to question meaning. No existing method achieves all three:
structural faithfulness, dynamic adaptation, and semantic conditioning at the constraint
level.

DCA-Trie fills this gap. Its TypeOracle conditions the valid token set on the knowledge
graph, the question semantics, and the partially generated chain. The constraint set
adapts at every decoding step, and only semantically relevant paths are admitted. The
result is a method that achieves structural faithfulness while reducing the search space
to paths the LLM actually needs to consider.
